\section{Collaborative Robotics in Industry 4.0}

The fourth industrial revolution, widely recognized as Industry 4.0, is fundamentally characterized by the integration of cyber-physical systems, the Industrial Internet of Things (IIoT), and advanced data analytics into manufacturing infrastructures. A cornerstone of this paradigm shift is the deployment of collaborative robotics. While early definitions of collaborative robots (cobots) focused strictly on physical safety for shared human-robot workspaces, modern industrial demands have expanded this definition to include complex, multi-agent automated systems \cite{lasi2014industry}.

The section below discusses the evolution of collaborative paradigms, the shift toward multi-robot systems, and the defining characteristics of cyber-physical assembly cells.

\subsection{From Human-Robot to Robot-Robot Collaboration}

Traditionally, collaborative robotics has been synonymous with Human-Robot Collaboration (HRC), where fences are removed to allow human operators and machines to work side-by-side \cite{matheson2019human}. However, in environments where tasks are fixed, highly repetitive, or require precision beyond human capability, the collaborative paradigm shifts toward Robot-Robot Collaboration (RRC) or Multi-Robot Collaboration (MRC). 

In RRC, the "collaboration" occurs between heterogeneous automated agents. Instead of a human handing a part to a machine, multiple robotic manipulators share a dynamic workspace to decompose tasks, execute parallel routines, and perform synchronized handovers. This approach eliminates human physical risk entirely while maximizing throughput, relying on the system's ability to seamlessly communicate state changes and trajectories between independent robot controllers \cite{kruger2009cooperation}.

\subsection{Characteristics of Multi-Robot Cells in Industry 4.0}

A multi-robot cell is classified as an Industry 4.0 system not by the presence of a human operator on the floor, but by its integration into a larger cyber-physical network. To achieve safe and efficient parallel execution without human intervention, these systems rely on several advanced architectures:

\begin{itemize}
    \item \textbf{Decentralized Communication:} Utilizing middleware, such as the Robot Operating System (ROS 2), to manage concurrent, multithreaded command streams across different hardware platforms (e.g., bridging an industrial manipulator with a research arm) \cite{macenski2022robot}.
    \item \textbf{Dynamic Collision Avoidance:} Implementing zone management and real-time spatial monitoring to prevent mechanical interference when multiple arms operate within intersecting workspaces \cite{khatib1986real,michalos2015design}.
    \item \textbf{Digital Twin Integration:} Providing a continuous, multi-mode virtual representation of the physical cell. This allows for real-time remote monitoring, trajectory teaching, and pre-execution validation before physical movement occurs \cite{glaessgen2012digital,wang2020digital,kritzinger2018digital}.
    \item \textbf{Remote Supervisory Control:} Shifting human interaction from physical proximity to remote, IIoT-enabled control rooms, elevating the human role to system oversight and high-level decision making \cite{wang2017collaborative}.
\end{itemize}

By relying on robust perception, strict safety logic, and virtual validation, modern multi-robot cells achieve the core goal of Industry 4.0: creating highly autonomous, resilient, and interconnected manufacturing environments \cite{lasi2014industry,wang2017collaborative}.
